% Canonical Paper IV section, adapted from the 2026-10-09 ported-AES candidate.
% Source/provenance and mathematical-status integration are recorded in governance.
The residual theory specifies when a state retains enough information for
future computation.  We now give a finite program interface that can be
attached to a punctured AES without identifying its geometry with its program.
Paper~0 owns the carrier geometry and its compactification
\cite{Yuan2026AEGPaperZero}; Paper~I supplies marked one-hole histories
\cite{Yuan2026AEGFoundations}.  The construction here is a structural proposal
for ported realizations.  Its finite-network and descent results are proved
with the hypotheses stated below.  The program/source distinction follows the
interface discipline of the Adva program--process core
\cite{Yuan2026AdvaProgramProcessCore}.

\subsection{Three kinds of opening}

An essential geometric puncture is a declared end where the given assignment
and metric do not extend as regular AES data.  A computational hole is an
ordered place for program substitution.  A port is an input or output interface
associated with a geometric end or a specified regular interface location.
Their identification requires additional data.

\begin{definition}[Ported AES realization proposal]
\label{def:ported-aes-realization}
Let $E=(M,g,a;\mu,\lambda)$ be an oriented regular AES with a declared
compactified carrier and marked end set $S$, where $|S|=k$.  A ported
realization consists of ordered typed port families $I,O$, an anchoring map
\[
  \pi:I\sqcup O\longrightarrow S\sqcup L,
\]
where $L$ is a declared set of regular interface locations, a finite typed
program decoration, and source, occurrence, boundary, and assembly records.
The map $\pi$ is not assumed injective or surjective.  Ports at a puncture carry
data in its end interface; they are not regular points with a value $a(p)$.
\end{definition}

The three integers $k,|I|,|O|$ are independent.  A three-output experiment
declares $|O|=3$; its internal input frontier may have any finite size.  If all
ports anchor injectively into $S$, then $k\ge |I|+|O|$; equality follows for a
bijection.  Ends without ports and ends carrying several ports are allowed.
The roles of the outputs must be declared.  They are not automatically three
scalar values or Adva's construction/space/time views; the outer interface and
the internal frontier are distinct \cite{Yuan2026AdvaTriadicFrontiers}.

Input binding retains the AES geometry unless a separate parameter family
$(g_\theta,a_\theta)$ is supplied.  Such a family needs a proof of the AES
identity, its admissible parameter domain, and its marked end behavior.
Binding is also distinct from topological capping.

\subsection{Finite programs, sources, and occurrences}

\begin{definition}[Strict finite arithmetic network]
\label{def:strict-finite-network}
An open network $P:A\rightsquigarrow B$ is a finite directed acyclic graph
with ordered typed input and output boundaries and ordered input/output slots
at every gate.  Each wire has one producer and one consumer, counting an
external output as a consumer.  Copy and discard are explicit gates.
Primitive gates denote deterministic partial functions.  Source identities are
immutable; distinct uses have distinct occurrence identities, even if their
sources or values agree.  External boundary positions likewise have occurrence
identities separate from source identities.
\end{definition}

The arithmetic fragment used below has exact rational constants, addition,
subtraction, multiplication, division, explicit copy, and strict discard.
Strict discard requires its input producer to be defined.  Thus
$\operatorname{discard}(1/0)$ remains outside the domain.  A finite acyclic
network cannot exhibit feedback divergence; feedback, side effects, and
unbounded interpretation need separately declared semantics.

Repeated use of a program template first creates separate named instances.
Their event occurrences are fresh; their original sources remain the same.
Administrative wire and event renaming then acts on those instances without
changing source or occurrence identity.  Here a source label names the declared
syntactic origin of a gate or wire occurrence.  The ancestry of a value supplied
after substitution is separately recovered from the connected producer DAG and
binding receipt; binding does not overwrite a body's immutable source labels.
Template reuse is therefore distinct from implicit sharing of one produced wire.

\subsection{Program substitution and its receipts}

Let $P$ have ordered holes $i_1,\ldots,i_m$, and let $q_1,\ldots,q_r$ be
already instantiated producer programs.  A binding $\beta$ matches producer
output occurrences to selected holes, preserving types and positions.  Every
supplied output is used once, and every selected hole is filled once.  Unused
outputs require an explicit discard; multiple uses require an explicit copy.
Holes not selected by a partial binding remain on the open frontier.  Fix the
external order as producer inputs in argument order, followed by the unfilled
body inputs.  Sequential partial filling with open producers can induce a
different order from simultaneous filling; comparison then requires a declared
typed boundary permutation.  The closed-producer example below requires none.

\begin{construction}[Certified simultaneous substitution]
\label{constr:certified-substitution}
Take disjoint administrative copies of the body and producer graphs, identify
the matched boundary wires, and retain all primitive gate occurrences and
sources.  Record the ordered binding, producer regions, body region, and call
nesting in a graft receipt.  The resulting network is
\[
  \Subst_\beta(P;\boldsymbol q).
\]
Accept it only when the resulting global graph is acyclic and its boundary
inventory is complete.  A numerical input is admitted by an explicit constant
producer with its source record, rather than by erasing the producing process.
\end{construction}

For ordinary one-way substitution, producers depend only on their external
inputs and the body depends on their outputs, so acyclicity follows.  General
seam assembly has no such guarantee: two locally acyclic pieces can form a
cycle after cross-wiring.  The notation $\Subst$ is reserved for program
substitution; $\Fill$ in Section~\ref{sec:rewrite-fibers} continues to denote
rewrite filling area.

\begin{proposition}[Type and identity preservation]
\label{prop:ported-substitution-preservation}
A legal substitution preserves boundary typing, primitive gate slot order,
source identities, and instantiated gate occurrences.  Distinct occurrences
are not identified by equality of numerical value.
\end{proposition}
\begin{proof}
The disjoint copies separate administrative names.  The binding changes only
which wire supplies a selected body input; it creates no primitive gate and
deletes none.  Its injectivity on producer output occurrences and selected
holes preserves the single-consumer rule.  Source and occurrence records are
transported by the stated renaming map.  The remaining boundaries retain their
declared order.  All other identifications are excluded by construction.
\end{proof}

\subsection{Substitution and partial evaluation}

Write $D_P$ for the strict numerical domain of $P$: inputs for which every
primitive gate required by the network is defined, including producers
eventually consumed by discard.  Primitive numerical functions are pure; no
primitive observes the scheduler.  The numerical result is therefore
independent of the chosen topological schedule.  An operational trace, its
first observed failure, and its costs can still depend on that schedule.

\begin{theorem}[Substitution--evaluation compatibility]
\label{thm:ported-substitution-evaluation}
Let $Q:X\rightsquigarrow A$ and $P:A\times Z\rightsquigarrow B$ be
compatible strict finite networks.  For the binding consuming all outputs of
$Q$, their composite has domain
\begin{equation}
  D_{P\circ Q}
  =\{(x,z):x\in D_Q,\ (\operatorname{Eval}(Q)(x),z)\in D_P\}.
  \label{eq:ported-composite-domain}
\end{equation}
On exactly this domain,
\begin{equation}
  \operatorname{Eval}(P\circ Q)(x,z)
  =\operatorname{Eval}(P)(\operatorname{Eval}(Q)(x),z).
  \label{eq:ported-composite-evaluation}
\end{equation}
The same statement applies to finitely many disjoint producers and partial
filling, retaining the unfilled boundary.
\end{theorem}
\begin{proof}
Use a topological order with producer gates before body gates.  A producer gate
has precisely the same ordered inputs, primitive partial function, and domain
predicate as in $Q$.  Induction gives the same values on its output wires
whenever the producer is defined.  The bound body inputs are therefore exactly
$\operatorname{Eval}(Q)(x)$.  A second topological induction through the body
gives its original gate values and domain predicates.  Every gate of the
composite is consequently defined if and only if both conditions in
Equation~\eqref{eq:ported-composite-domain} hold.  This proves both equations.
Purity permits a different topological order without changing the successful
result or whether all gates are defined; it does not identify ordered traces.
\end{proof}

\begin{corollary}[Three equality levels for composition]
\label{cor:ported-composition-equalities}
Compatible composition has the following laws.
\begin{enumerate}
  \item Partial numerical functions satisfy identity and associativity
  strictly, including equality of their domains.
  \item Operation networks satisfy these laws up to an explicit isomorphism
  of administrative names and contraction of administrative port segments.
  The isomorphism preserves instantiated occurrences, sources, types, slot
  orders, and ordered external boundaries.  Identity networks use the same
  source and occurrence contract as the interface object they preserve.
  \item Literal receipts for different calls or parenthesizations remain
  distinct.  Their final operation networks are related by the preceding
  transport certificate, not by deleting the call history.
\end{enumerate}
\end{corollary}
\begin{proof}
Equation~\eqref{eq:ported-composite-domain} applied twice gives the same three
domain conditions for either parenthesization;
Equation~\eqref{eq:ported-composite-evaluation} then gives the same result.
At graph level, both constructions contain the same instantiated primitive
gates and the same final boundary identifications.  Map each renamed gate and
wire to its counterpart and contract only administrative segments.  This
supplies the stated isomorphism.  A receipt records the actual calls and their
nesting, which these operations do not make equal.
\end{proof}

\subsection{Joint geometric and computational descent}

\begin{proposition}[Collar and interface descent]
\label{prop:ported-collar-descent}
Suppose finitely many oriented regular AES pieces have the same signed
parameters $(\mu,\lambda)$, and orientation-compatible open collar maps
$\phi_{vw}$ produce a Hausdorff smooth surface.  Require inverse and cocycle
conditions and, on every overlap,
\[
  a_v=a_w\circ\phi_{vw},\qquad g_v=\phi_{vw}^{*}g_w.
\]
Require computational seam records to agree on ordered ports, types, sources,
occurrences, gate parameters, and boundary wire identities, with a complete
assembly record for nonlocal relations.  Require the assembled global program
to be a finite DAG.  Require marked end identities and the anchoring map $\pi$
to be compatible with all identifications.

Then the AES data and the computational network descend together.  The
canonical arithmetic frames transform by the collar Jacobians.  The assembled
program obeys Theorem~\ref{thm:ported-substitution-evaluation} for every
compatible causal composition across its interfaces.
\end{proposition}
\begin{proof}
Overlap equality descends the smooth function and positive definite metric.
The identity $|\mathrm da|_g^2=\mu^2+\lambda^2a^2$ is tensorial and holds in
every local chart, hence globally.  Compatible orientation and equal signed
parameters give the same canonical frame formula in each chart, so those
frames also descend.  The seam inventory identifies precisely the specified
wire occurrences without merging sources or introducing implicit copies.
Compatible end markings and anchoring maps descend the attachment of ports to
the carrier.  Global acyclicity yields a topological execution order.  The
numerical composition result follows from the finite-network theorem.
\end{proof}

\begin{warning}[The motion bridge remains additional]
\label{warn:ported-motion-bridge}
Proposition~\ref{prop:ported-collar-descent} does not prove that a gate of the
attached network is an AES motion.  Such a realization requires primitive
partial state transformations $\Phi_\omega$, exact domains, parameter/source
transport, and overlap equations
\[
  T_{vw}\circ\Phi^v_\omega=\Phi^w_\omega\circ T_{vw}
\]
on their stated domains.  Binary gates, explicit copy/discard, and negative
scalings cannot be inferred from the single-value positive affine model.
Full-state tangential residuals must also be retained: an identity arithmetic
word on assignment values need not be the identity on AES states.  Attaching a
program to a surface does not establish a faithful arithmetic-motion
realization.
\end{warning}

For an already given genus-zero $k$-end carrier, $k\ge3$, a pants
decomposition has $k-2$ three-end pieces and $k-3$ internal seams.  Restrict
all data and retain the overlap and assembly records.  Joint descent then
recovers the original ported object.  A second decomposition gives another
atlas of that same object.  Each three-end piece may have punctures and seam
boundaries; it is not thereby a complete three-cusp AES.  Conformal point
compactification additionally requires punctured-disc ends.  These geometric
constructions and distinctions belong to Paper~0
\cite{Yuan2026AEGPaperZero}.

Local block frontiers can alternate across a seam.  Even when the global event
graph is acyclic, the coarse graph whose vertices are whole geometric pieces
may contain a cycle.  One then cannot execute each whole piece just once in a
block order; execution uses the global causal event order and the exposed
seam frontier.

\subsection{Features and future computation}

Fix the experiment $(P_C,M_C,O_C)$ of
Section~\ref{sec:contextual-residuals}, with its declared treatment of illegal
arithmetic.  Use a domain-compositional partial action, or the free
continuation word monoid, before totalization; do not import unproved operation
relations.  Divergence is not identified with a type error or division by zero.
Here $\cdot$ is a declared total right action on $P_C^\bot$, for example
execution of continuation words with an absorbing illegal state and a distinct
invalid observation.  Contextual equivalence means
\[
  p\equiv_Cp'
  \quad\Longleftrightarrow\quad
  O_C(p\cdot\eta)=O_C(p'\cdot\eta)
  \quad\text{for every }\eta\in M_C.
\]
A right congruence is an equivalence relation preserved by every continuation.

Let $\chi_C:P_C^\bot\to Z_C$ be a proposed feature.  It may retain a scalar,
operator, marking, source/occurrence package, or residual.  The kind of feature
and its observation policy are part of the interface.

\begin{theorem}[Feature observation descent]
\label{thm:feature-observation-descent}
The following conditions are equivalent.
\begin{enumerate}
  \item For every declared continuation $\eta$, there is a function $o_\eta$
  on the reachable feature image satisfying
  \[
    o_\eta(\chi_C(p))=O_C(p\cdot\eta).
  \]
  \item $\ker\chi_C$ is contained in contextual continuation equivalence:
  \[
    \chi_C(p)=\chi_C(p')\quad\Longrightarrow\quad p\equiv_Cp'.
  \]
\end{enumerate}
Legal-domain differences are included through the declared invalid outcome.
\end{theorem}
\begin{proof}
The first condition makes all future observations equal on each feature fiber,
which is exactly the second condition.  Conversely, define $o_\eta$ using any
representative of a feature fiber; the second condition makes the choice
immaterial.  This is existence of observation factorization, not an assertion
that the factors can be computed efficiently or at all.
\end{proof}

\begin{proposition}[Online feature update]
\label{prop:feature-online-update}
Feature transitions $u_\eta$ satisfying
\[
  u_\eta(\chi_C(p))=\chi_C(p\cdot\eta)
\]
exist on the reachable image if and only if $\ker\chi_C$ is a right
congruence.  Together with factorization of $O_C$, they give an exact feature
continuation machine.  An executable implementation still requires effective
transition and observation maps.
\end{proposition}
\begin{proof}
Necessity follows by applying $u_\eta$ to equal feature values.  For sufficiency,
define $u_\eta$ by any representative.  Right congruence makes it well defined,
and the action laws descend from the given action.  Output factorization
supplies the current observation.
\end{proof}

\begin{example}[Observation sufficiency without online update]
\label{ex:feature-not-online}
Take states $0,1,2$, constant observation zero, and the idempotent update
$f(0)=0$, $f(1)=2$, $f(2)=2$.  Every state is contextually equivalent for the
monoid $\{\id,f\}$.  Let $\chi(0)=\chi(1)=0$, $\chi(2)=1$.
Feature observation descent holds, but $\chi(f(0))\ne\chi(f(1))$, so no
feature update can be defined on input feature zero.  A future-sufficient
observation and a closed online state are different requirements.
\end{example}

Actual opposite-edge characteristic extraction remains a separate witness
problem \cite{Yuan2026AdvaArithmeticCalibration}.  If a spectral feature is
used, first declare the endomorphism closure, operator domain, marking, and
residual.  A pairwise relation does not automatically have an ordinary
spectrum.  Three-side feedback also requires a snapshot and update schedule or
a separately proved fixed-point semantics.

\subsection{A four-puncture, four-input, three-output experiment}

\begin{example}[Fixed AES carrier with program interfaces]
\label{ex:ported-four-puncture}
On $M=\PP^1(\C)\setminus\{-1,0,1,\infty\}$, put
\begin{align*}
  a&=\frac{x^4}{4}-\frac{x^2}{2}+y^2,
  & A&=x^3-x,& B&=2y,\\
  L&=A^2+B^2,& F&=1+a^2,& g&=\frac{L}{F}|\mathrm dz|^2.
\end{align*}
Off the three finite deleted points, $L>0$ and $|\mathrm da|_g^2=F$.
At each finite point $\mathrm da=0$, excluding a nondegenerate regular AES
extension.  Along the imaginary axis $a\to\infty$, excluding a real-valued
extension at infinity.  Thus all four marked ends are essential.  This direct
verification does not use a general puncture/critical-point assertion.

Anchor the four input holes to $-1,0,1,\infty$, in that order, keeping input
values as port data.  The finite assignment limits $(-1/4,0,-1/4)$ and the divergent infinite limit
are not replaced by those values.  Anchor the three output ports at $i,1+i,-1+i$, respectively; their output
values likewise need not equal the assignment values at those regular sites.  Choose producers
\[
  q_1=1+1=2,\quad q_2=3+4=7,\quad
  q_3=6/2=3,\quad q_4=9-4=5.
\]
The body, with an explicit copy for each twice-used input, reads out
\[
  y_1=x_1+x_2,\qquad y_2=x_3x_4,\qquad
  y_3=\frac{x_1+x_3}{x_2-x_4}.
\]
The output is exactly $(9,15,5/2)$.  It is a declared scalar projection of
this experiment, not an identification with Adva's native three output roles.
The run retains eight constant sources, their occurrences and explicit copy
records, the operation DAG, and the graft receipt.
\end{example}

\begin{figure}[ht]
\centering
\begin{tikzpicture}[
  node distance=8mm and 15mm,
  box/.style={draw,rounded corners,align=center,minimum height=24mm},
  arr/.style={-{Latex[length=2mm]}}
]
\node[align=right] (inputs) {$q_1\mapsto2\quad h_1$\\[2mm]
$q_2\mapsto7\quad h_2$\\[2mm]$q_3\mapsto3\quad h_3$\\[2mm]
$q_4\mapsto5\quad h_4$};
\node[box,right=of inputs] (body) {typed body\\explicit copy\\strict division\\source ledger};
\node[align=left,right=of body] (outputs) {$y_1=9$\\[5mm]$y_2=15$\\[5mm]$y_3=5/2$};
\draw[arr] (inputs) -- (body);
\draw[arr] (body) -- (outputs);
\end{tikzpicture}
\caption{A fixed four-end AES carrier with four computational inputs and
three declared scalar outputs.  The diagram displays the program interface,
not paths of arithmetic gates through the AES.}
\label{fig:ported-four-inputs}
\end{figure}

Use either cut $|z-c|=3/4$ with $c=-1/2$ or $c=1/2$.  The inner and outer
charts on the overlap are
\[
  \xi=\frac{z-c}{3/4},\qquad
  \eta=\frac{3/4}{z-c},\qquad \eta=1/\xi.
\]
Each cut separates two marked ends from the other two.  For each cut, also
declare an allocation of gate occurrences to local packets; this allocation is
additional program data, not a consequence of the port anchoring map.
Restrict the carrier data to the chosen pieces and retain each local event
region, its exposed wire frontier, and the complete ordered seam/assembly
contract.
Reciprocal coordinate transport is an exact change of chart.  The tensor and
chain-rule argument in Proposition~\ref{prop:ported-collar-descent} recovers
the same global AES and operation network from either decomposition.  Local
seam source references identify transported data; they are not new constant
sources.  Execution after assembly follows the global event order, which need
not be a whole-piece block order.  No path of a binary gate through the AES
is asserted.

\subsection{Finite checks and remaining obligations}

The accompanying independent finite checker uses exact rational arithmetic for
this example and explicit positive and negative controls.  Its scope and
reproduction commands are recorded in the Paper~IV README.  It is a new
verification of the integrated model, not a rerun of the unavailable checker
reported with the candidate chapter.  Finite checks supplement the proofs;
they do not prove continuous collar identities or a faithful motion bridge.

The important negative controls are implicit sharing, duplicate binding,
mismatched types, unaccounted seam data, and cycles formed by locally acyclic
pieces.  A literal constant two and the producer $1+1$ have the same scalar
output but different sources and operation histories.  Strict evaluation
retains failure of a discarded $1/0$ producer.  Different association receipts
remain distinct, and Example~\ref{ex:feature-not-online} remains an obstruction
to online feature update.

The remaining obligations are stated as
Problems~\ref{prob:faithful-multiwire-aeg} and
\ref{prob:opposite-edge-features}: realize the actual typed generators by AES
state transformations and construct the actual opposite-edge feature
witnesses.  Neither is a consequence of finite program substitution.
